\begin{thebibliography}{00}

\bibitem{ref1}
A. T. D. Perera, K. Javanroodi et al., ``Challenges resulting from urban density and climate change for the EU energy transition,'' \emph{Nature Energy}, vol.~8, no.~4, pp.~397--412, 2023.

\bibitem{ref2}
R. van Heerden, O. Y. Edelenbosch et al., ``Demand-side strategies enable rapid and deep cuts in buildings and transport emissions to 2050,'' \emph{Nature Energy}, vol.~10, no.~3, pp.~380--394, 2025.

\bibitem{ref3}
B. Zhou and M. Hu, ``Impact of architectural design elements on the energy efficiency and thermal comfort in neighborhood building ensembles,'' \emph{Journal of Building Engineering}, vol.~100, 2025.

\bibitem{ref4}
L. Du, H. Wang et al., ``Impact of block form on building energy consumption, urban microclimate and solar potential: A case study of Wuhan, China,'' \emph{Energy and Buildings}, vol.~328, 2025.

\bibitem{ref5}
B. Liu, Y. Liu et al., ``Urban morphology indicators and solar radiation acquisition: 2011--2022 review,'' \emph{Renewable and Sustainable Energy Reviews}, vol.~199, 2024.

\bibitem{ref6}
K. Javanroodi, A. T. D. Perera et al., ``Designing climate resilient energy systems in complex urban areas considering urban morphology: A technical review,'' \emph{Advances in Applied Energy}, vol.~12, 2023.

\bibitem{ref7}
A. T. D. Perera, K. Javanroodi, and V. M. Nik, ``Climate resilient interconnected infrastructure: Co-optimization of energy systems and urban morphology,'' \emph{Applied Energy}, vol.~285, 2021.

\bibitem{ref8}
Z. Shi, J. A. Fonseca, and A. Schlueter, ``A review of simulation-based urban form generation and optimization for energy-driven urban design,'' \emph{Building and Environment}, vol.~121, pp.~119--129, 2017.

\bibitem{ref9}
A. Eltaweel and Y. Su, ``Parametric design and daylighting: A literature review,'' \emph{Renewable and Sustainable Energy Reviews}, vol.~73, pp.~1086--1103, 2017.

\bibitem{ref10}
S. Longo, F. Montana, and E. Riva Sanseverino, ``A review on optimization and cost-optimal methodologies in low-energy buildings design and environmental considerations,'' \emph{Sustainable Cities and Society}, vol.~45, pp.~87--104, 2019.

\bibitem{ref11}
F. Ascione, R. F. De Masi et al., ``Optimization of building envelope design for nZEBs in Mediterranean climate: Performance analysis of a residential case study,'' \emph{Applied Energy}, vol.~183, pp.~938--957, 2016.

\bibitem{ref12}
G. Li, Q. He et al., ``Accelerated inverse urban design: A multi-objective optimization method to photovoltaic power generation potential, environmental performance and economic performance in urban blocks,'' \emph{Sustainable Cities and Society}, vol.~120, 2025.

\bibitem{ref13}
K. Liu, X. Xu et al., ``A multi-objective optimization framework for designing urban block forms considering daylight, energy consumption, and photovoltaic energy potential,'' \emph{Building and Environment}, vol.~242, 2023.

\bibitem{ref14}
A.-T. Nguyen, S. Reiter, and P. Rigo, ``A review on simulation-based optimization methods applied to building performance analysis,'' \emph{Applied Energy}, vol.~113, pp.~1043--1058, 2014.

\bibitem{ref15}
V. Machairas, A. Tsangrassoulis, and K. Axarli, ``Algorithms for optimization of building design: A review,'' \emph{Renewable and Sustainable Energy Reviews}, vol.~31, pp.~101--112, 2014.

\bibitem{ref16}
Y. Zhang, B. K. Teoh, and L. Zhang, ``Multi-objective optimization for energy-efficient building design considering urban heat island effects,'' \emph{Applied Energy}, vol.~376, 2024.

\bibitem{ref17}
G. Shi, S. Yao et al., ``Multi-performance collaborative optimization of existing residential building retrofitting in extremely arid and hot climate zones: A case study in Turpan, China,'' \emph{Journal of Building Engineering}, vol.~89, 2024.

\bibitem{ref18}
Y. Pan, M. Zhu et al., ``Building energy simulation and its application for building performance optimization: A review of methods, tools, and case studies,'' \emph{Advances in Applied Energy}, vol.~10, 2023.

\bibitem{ref19}
Z. Zhao, H. Li, and S. Wang, ``Surrogate-assisted coordinated design optimization of building and microclimate considering their mutual impacts,'' \emph{Applied Energy}, vol.~383, 2025.

\bibitem{ref20}
P. Westermann and R. Evins, ``Surrogate modelling for sustainable building design: A review,'' \emph{Energy and Buildings}, vol.~198, pp.~170--186, 2019.

\bibitem{ref21}
K. Li, T. Zhang, and R. Wang, ``Deep Reinforcement Learning for Multiobjective Optimization,'' \emph{IEEE Transactions on Cybernetics}, vol.~51, no.~6, pp.~3103--3114, 2021.

\bibitem{ref22}
M. Khoroshiltseva, D. Slanzi, and I. Poli, ``A Pareto-based multi-objective optimization algorithm to design energy-efficient shading devices,'' \emph{Applied Energy}, vol.~184, pp.~1400--1410, 2016.

\bibitem{ref23}
R. Liu, T. Fang et al., ``Controllable cross-building multi-objective optimisation for NZEBs: A framework utilising parametric generation and intelligent algorithms,'' \emph{Applied Energy}, vol.~374, 2024.

\bibitem{ref24}
S. Miraba, A. Salehi et al., ``Adaptive BIPV shading optimization for electricity generation and building electricity management, sDA, and DGP using multi-objective algorithms,'' \emph{Applied Energy}, vol.~402, 2025.

\bibitem{ref25}
A. Ciardiello, F. Rosso et al., ``Multi-objective approach to the optimization of shape and envelope in building energy design,'' \emph{Applied Energy}, vol.~280, 2020.

\bibitem{ref26}
X. Wang, S. Wang et al., ``Deep Reinforcement Learning: A Survey,'' \emph{IEEE Transactions on Neural Networks and Learning Systems}, vol.~35, no.~4, pp.~5064--5078, 2024.

\bibitem{ref27}
H. Jemin, L. Joonho et al., ``Learning agile and dynamic motor skills for legged robots,'' \emph{Science Robotics}, vol.~4, no.~26, 2019.

\bibitem{ref28}
O. Vinyals, I. Babuschkin et al., ``Grandmaster level in StarCraft II using multi-agent reinforcement learning,'' \emph{Nature}, vol.~575, no.~7782, pp.~350--354, 2019.

\bibitem{ref29}
X. Wang, W. Chen et al., ``Video Captioning via Hierarchical Reinforcement Learning,'' in \emph{2018 IEEE/CVF Conference on Computer Vision and Pattern Recognition (CVPR)}, 2018, pp.~4213--4222.

\bibitem{ref30}
J. Li, L. Yao et al., ``Deep reinforcement learning for pedestrian collision avoidance and human-machine cooperative driving,'' \emph{Information Sciences}, vol.~532, pp.~110--124, 2020.

\bibitem{ref31}
V. Mnih, K. Kavukcuoglu et al., ``Human-level control through deep reinforcement learning,'' \emph{Nature}, vol.~518, no.~7540, pp.~529--533, 2015.

\bibitem{ref32}
L. Timothy P, H. Jonathan J et al., ``Continuous control with deep reinforcement learning,'' \emph{arXiv preprint arXiv:1509.02971}, 2015.

\bibitem{ref33}
Y. Du, H. Zandi et al., ``Intelligent multi-zone residential HVAC control strategy based on deep reinforcement learning,'' \emph{Applied Energy}, vol.~281, 2021.

\bibitem{ref34}
S. Li, S. Su, and X. Lin, ``Optimizing the hyper-parameters of deep reinforcement learning for building control,'' \emph{Building Simulation}, vol.~18, no.~4, pp.~765--789, 2025.

\bibitem{ref35}
Z. Wu, Y. Li et al., ``Distributed voltage control for multi-feeder distribution networks considering transmission network voltage fluctuation based on robust deep reinforcement learning,'' \emph{Applied Energy}, vol.~379, 2025.

\bibitem{ref36}
P. Kou, D. Liang et al., ``Safe deep reinforcement learning-based constrained optimal control scheme for active distribution networks,'' \emph{Applied Energy}, vol.~264, 2020.

\end{thebibliography}
